\section{三角関数と漸化式}\label{節：三角関数と漸化式}
\subsection{一定角を足す場合}

単位円上の点を$z_n=e^{i\theta_n}$と書き,\,一定角$\phi$ずつ進めると
\[
 z_{n+1}=e^{i\phi}z_n,\qquad
 \theta_n=\theta_1+(n-1)\phi
\]
となる.
横座標$a_n=\cos\theta_n$だけを記録すると,\,加法公式から
\[
 \cos(\theta_n+2\phi)+\cos\theta_n
 =2\cos\phi\cos(\theta_n+\phi)
\]
なので
\[
 a_{n+2}=2\cos\phi\,a_{n+1}-a_n.
\]
一つの横座標だけでは点の上下を区別できないが,\,隣り合う二つの記録があれば次を更新できる.
特性方程式の根は$e^{i\phi},e^{-i\phi}$であり,\,前節につながる.
$0<\phi<\pi$では二根が非実数となる.

一方,\,傾き$a_n=\tan\theta_n$を記録すれば
\[
 a_{n+1}=\frac{a_n+\tan\phi}{1-a_n\tan\phi}
\]
という1次分数型になる.
ただし$\tan\phi$と各項の正接が定義され,\,更新の分母が非零の範囲で考える.
同じ角度の動きでも,\,\textbf{横座標と傾きでは更新式の形が変わる}.

\subsection{角度を倍にする場合}

次は
\[
 a_1=\frac{\sqrt3}{2},\qquad a_{n+1}=2a_n^2-1
\]
を考える.
倍角公式$\cos2\theta=2\cos^2\theta-1$と,
$a_1=\cos\frac{\pi}{6}$を使えば,
\[
 a_n=\cos\left(2^{n-1}\frac{\pi}{6}\right)
\]
と予想できる.
初項に合い,\,この式を更新の右辺に代入すると次の項の式になるので,\,帰納法で正しい.

角度は$\theta_{n+1}=2\theta_n$と選べる.
複素数では
\[
 z_{n+1}=z_n^2,\qquad z_n=z_1^{\,2^{n-1}}
\]
であり,$a_n=\operatorname{Re}z_n$である.
\NewFigDoubling

これは一定角の回転ではない.
現在の角度$\theta_n$をもう一度足すため,\,毎回の回転角が変わる.
さらに$z$と$-z$は同じ$z^2$へ移るので,\,円周全体の写像として一対一ではない.

余弦の値から角度は一意に決まらない.
そのため$a_{n+1}=2a_n^2-1$から,\,選び方によらず$\theta_{n+1}=2\theta_n$が強制されると読むのは正確でない.
整合する角度の列を一つ選び,\,そこから作った候補が元の式を満たすことを示している.

\subsection{正接の倍角と定義域}

傾きを記録する倍角型は
\[
 a_{n+1}=\frac{2a_n}{1-a_n^2}
\]
である.
$a_n=\tan\theta_n$と表せれば,$\theta_{n+1}=2\theta_n$と選べる.
ただし$a_n=\pm1$では次の傾きが縦向きになり,\,実数の値としては定義できない.
円周上の動きが続いても,\,傾きの列が続くとは限らない.

\noindent\textbf{【例題】}\par\nobreak
\[
 a_1=2-\sqrt3,\qquad
 a_{n+1}=\frac{2a_n}{1-a_n^2}
\]
の一般項を求め,\,全項が定義されることを示せ.

\noindent\textbf{【方針】}\par\nobreak
$2-\sqrt3=\tan\frac{\pi}{12}$です. 角度を倍にし,\,途中で縦向きにならないか確認します.

\noindent\textbf{【解答】}\par\nobreak
\[
 a_n=\tan\left(2^{n-1}\frac{\pi}{12}\right).
\]
\noindent\textbf{【解法】}\par\nobreak
初項と倍角公式から候補を得る.
$n=1$では正接は定義できる.
$n\ge2$でその角度が$\frac{\pi}{2}+\ell\pi$に一致するとすれば
\[
 2^{n-2}=3+6\ell
\]
となるが,2の累乗は3の倍数でないので不可能である.
従って全項の正接が定義される.
また次の角度も縦向きにならないため,\,更新の分母$1-a_n^2$も非零である.
帰納法により一般項が成立する.

なお,$\Phi$は前節の実軸と円の対応であり,
\[
 \Phi\left(\frac{2a}{1-a^2}\right)=\Phi(a)^2
\]
となる.
正接の倍角型も,\,円周上の二乗を別の座標で記録していた.
1次分数型の一定角回転とは,\,同じ単位円を使っても更新の規則が違う.

\subsection{整数倍角とチェビシェフ多項式}

角度を2倍する式を見たので,\,一般の整数倍も考えよう.
多項式$T_m(x)$を
\[
 T_0(x)=1,\quad T_1(x)=x,
\]
\[
 T_{m+1}(x)=2xT_m(x)-T_{m-1}(x)
\]
で定める. これを\textbf{チェビシェフ多項式}という.
加法公式によって
\[
 T_m(\cos\theta)=\cos(m\theta)
\]
が帰納的に示される.
例えば
\[
 T_2(x)=2x^2-1,\qquad T_3(x)=4x^3-3x.
\]
特定の$x$を固定すると,$m$に関しては2階の線形漸化式である.

一方,\,固定した整数$d\ge2$について
\[
 a_{n+1}=T_d(a_n)
\]
を反復すると,$a_1=\cos\theta$のとき
\[
 a_n=\cos(d^{n-1}\theta)
\]
となる. こちらは$n$に関して非線形の更新である.
多項式の次数を進める添字$m$と,\,更新回数を表す$n$を区別しよう.

さらに
\[
 T_j(T_m(x))=T_{jm}(x)
\]
という合成の公式もある.
$x=\cos\theta$で確かめると両辺は$\cos(jm\theta)$になる.
両辺は多項式で,\,区間$[-1,1]$の無数の点で一致するため,\,すべての実数・複素数$x$で一致する.
零でない多項式の根が次数を超えないことを使っている.
従って一般の初項でも
\[
 a_n=T_{d^{n-1}}(a_1)
\]
と表せる.

\subsection{三角関数の外側にも続く式}

実数$a_1$が$[-1,1]$の外にあると,\,実角の余弦としては表せない.
しかし多項式の更新自体は定義されている.
任意の$x$に対し,\,非零の複素数$u$を
\[
 u^2-2xu+1=0
\]
の根として選べば,$x=\frac{u+u^{-1}}2$であり,
\[
 T_m(x)=\frac{u^m+u^{-m}}2
\]
となる. 初めの二つの値と更新を確かめればよい.
$|u|=1$なら余弦の表示に戻り,\,実数$x>1$なら実根を選んで指数的な増大として読める.

例えば$a_1=\frac54$,$a_{n+1}=2a_n^2-1$では$u=2$と選べるので,
\[
 a_n=\frac{2^{\,2^{n-1}}+2^{-2^{n-1}}}{2}.
\]
三角関数の公式は,\,二つの逆数の冪を足した式の一つの姿でもあった.
第2章の等比的な成分への分解と,\,ここでも再会する.
